\chapter*{Résumé}

Ce rapport présente le projet de stage \textbf{SmartTrans-AI}, plateforme de gestion de transport enrichie par l'intelligence artificielle. L'application permet de planifier et de suivre des missions, de gérer les clients, les chauffeurs et les véhicules, et d'alerter sur les échéances critiques (permis de conduire, assurance, visite technique). Trois profils sont distingués : administrateur, responsable transport et chauffeur. L'authentification repose sur JSON Web Token (JWT) et des rôles. Un journal d'activités, en lecture seule, trace automatiquement les opérations. Un assistant conversationnel, ancré sur les données métier, répond aux questions opérationnelles et propose des ressources disponibles. La réalisation s'appuie sur Spring Boot (API REST), PostgreSQL, React et Vite. Ce document décrit le cadre du stage, l'analyse et la conception (cas d'utilisation, classes), la réalisation, les tests et les perspectives.

\vspace{0.8em}
\noindent\textbf{Mots-clés :} gestion de transport, missions, JWT, Spring Boot, React, assistant IA, audit.
